\documentclass[10pt,a4paper]{amsart}
\usepackage{amsmath,amssymb,amsthm,graphicx}
% For commutative diagrams we use Xy-pic
\usepackage[all]{xy}

% We use 2cell for 2-commutative diagrams.
\xyoption{2cell}
% Theorem environments.
%
\theoremstyle{plain}
\newtheorem{theorem}[subsection]{Theorem}
\newtheorem{proposition}[subsection]{Proposition}
\newtheorem{lemma}[subsection]{Lemma}

\theoremstyle{definition}
\newtheorem{definition}[subsection]{Definition}
\newtheorem{example}[subsection]{Example}
\newtheorem{exercise}[subsection]{Exercise}
\newtheorem{situation}[subsection]{Situation}

\theoremstyle{remark}
\newtheorem{remark}[subsection]{Remark}
\newtheorem{remarks}[subsection]{Remarks}

\numberwithin{equation}{subsection}

% Macros
%
\def\lim{\mathop{\mathrm{lim}}\nolimits}
\def\colim{\mathop{\mathrm{colim}}\nolimits}
\def\Spec{\mathop{\mathrm{Spec}}}
\def\Hom{\mathop{\mathrm{Hom}}\nolimits}
\def\Ext{\mathop{\mathrm{Ext}}\nolimits}
\def\SheafHom{\mathop{\mathcal{H}\!\mathit{om}}\nolimits}
\def\SheafExt{\mathop{\mathcal{E}\!\mathit{xt}}\nolimits}
\def\Sch{\mathit{Sch}}
\def\Mor{\mathop{\mathrm{Mor}}\nolimits}
\def\Ob{\mathop{\mathrm{Ob}}\nolimits}
\def\Sh{\mathop{\mathit{Sh}}\nolimits}
\def\NL{\mathop{N\!L}\nolimits}
\def\CH{\mathop{\mathrm{CH}}\nolimits}
\def\proetale{{pro\text{-}\acute{e}tale}}
\def\etale{{\acute{e}tale}}
\def\QCoh{\mathit{QCoh}}
\def\Ker{\mathop{\mathrm{Ker}}}
\def\Im{\mathop{\mathrm{Im}}}
\def\Coker{\mathop{\mathrm{Coker}}}
\def\Coim{\mathop{\mathrm{Coim}}}

% Boxtimes
%
\DeclareMathSymbol{\boxtimes}{\mathbin}{AMSa}{"02}

%
% Macros for moduli stacks/spaces
%
\def\QCohstack{\mathcal{QC}\!\mathit{oh}}
\def\Cohstack{\mathcal{C}\!\mathit{oh}}
\def\Spacesstack{\mathcal{S}\!\mathit{paces}}
\def\Quotfunctor{\mathrm{Quot}}
\def\Hilbfunctor{\mathrm{Hilb}}
\def\Curvesstack{\mathcal{C}\!\mathit{urves}}
\def\Polarizedstack{\mathcal{P}\!\mathit{olarized}}
\def\Complexesstack{\mathcal{C}\!\mathit{omplexes}}
% \Pic is the operator that assigns to X its picard group, usage \Pic(X)
% \Picardstack_{X/B} denotes the Picard stack of X over B
% \Picardfunctor_{X/B} denotes the Picard functor of X over B
\def\Pic{\mathop{\mathrm{Pic}}\nolimits}
\def\Picardstack{\mathcal{P}\!\mathit{ic}}
\def\Picardfunctor{\mathrm{Pic}}
\def\Deformationcategory{\mathcal{D}\!\mathit{ef}}
\newenvironment{reference}{}{}
\newenvironment{slogan}{}{}
\newenvironment{history}{}{}
\title{A symmetric recurring real matrix is negative semidefinite, with kernel exactly the line spanned by its positive vector}
\author{The Stacks Project Authors}
\date{stacks-project:models.tex:local-snapshot-20261012; GNU FDL 1.2 or later}
\begin{document}
\maketitle
\noindent\textit{Editorial note.} The counted unit is the Stacks Project Lemma lemma-recurring-symmetric-real at source lines 372-385 together with its proof at source lines 387-405, a straight adjacent pairing in the chapter file. The chapter is the upstream file models.tex, which is the file name recorded in the source record; the pool directory carries the longer name stacks-models and is not used as the chapter name anywhere in this card.\par\medskip
\noindent\textit{Editorial note.} The author himself heads the retained proof environment with the words First proof, an optional argument of the proof environment, which states that a second argument for the same lemma is to be found further on in the chapter. Only the first proof is retained here; the counted unit is the lemma with the argument that the author places next to it, and the wording of the header is kept verbatim in the retained span.\par\medskip
\noindent\textit{Editorial note.} This chapter's candidate index records the statement body and the proof body without their closing markers, so the retained spans used here are the full environments: the statement span ends at the line of the closing marker of the lemma environment and the proof span starts at the opening marker of the proof environment. The same one-line accounting was recorded in the earlier card of this chapter, P2-models-recurring-real.\par\medskip
\noindent\textit{Editorial note.} Neither the retained statement nor the retained proof contains a reference to another labelled result of the chapter and neither contains a citation, so this card carries no navigation pointer and no bibliography entry.\par\medskip
\noindent\textit{Editorial note.} The retained preamble is the chapter preamble's Xy-pic setup at lines 14-18, the theorem declarations and mathematical macros at lines 154-217, and the three no-op wrapper environments.\par\medskip
\noindent\textit{Editorial note.} The proof is the author's own complete argument and contains no omission wording: the rescaling, the regrouping identity for the sum of squared coordinate differences and the equality case are all written out.\par\medskip
\noindent\textit{Allowed standard prerequisites (editorial).}\begin{itemize}
\item Recurring matrices and the quadratic form attached to them. A matrix whose off-diagonal entries are nonnegative together with a positive vector in its kernel, with no proper nonempty subset of indices that fails to interact with its complement, is what the chapter calls a symmetric recurring matrix; the quadratic form of such a matrix is written with a transposed vector on the left, and the surrounding lemmas of the same chapter use the recurrence hypothesis to control the kernel and the rank. The author's own hypotheses are retained inside the counted statement.. The counted lemma is the semidefiniteness statement that the neighbouring lemmas of the chapter invoke when they compute the rank of a recurring matrix and the kernel of the matrices attached to a numerical type.
\end{itemize}\medskip
\begin{lemma}
\label{lemma-recurring-symmetric-real}
Let $A = (a_{ij})$ be a symmetric real $n \times n$ matrix with
$a_{ij} \geq 0$ for $i \not = j$.
Let $m = (m_1, \ldots, m_n)$ be a real vector with $m_i > 0$.
Assume
\begin{enumerate}
\item $Am = 0$,
\item there is no proper nonempty subset $I \subset \{1, \ldots, n\}$
such that $a_{ij} = 0$ for $i \in I$ and $j \not \in I$.
\end{enumerate}
Then $x^t A x \leq 0$ with equality if and only if $x = qm$
for some $q \in \mathbf{R}$.
\end{lemma}

\begin{proof}[First proof]
After replacing $a_{ij}$ by $a_{ij}m_im_j$ we may assume $m_i = 1$
for all $i$. Condition (1) means $-a_{ii} = \sum_{j \not = i} a_{ij}$
for all $i$. Recall that $x^tAx = \sum_{i, j} x_ia_{ij}x_j$.
Then
\begin{align*}
\sum\nolimits_{i \not = j} -a_{ij}(x_j - x_i)^2 & =
\sum\nolimits_{i \not = j} -a_{ij}x_j^2 + 2a_{ij}x_ix_i - a_{ij}x_i^2 \\
& =
\sum\nolimits_j a_{jj} x_j^2 +
\sum\nolimits_{i \not = j} 2a_{ij}x_ix_i +
\sum\nolimits_j a_{jj} x_i^2 \\
& = 2x^tAx
\end{align*}
This is clearly $\leq 0$. If equality holds, then let $I$ be the set
of indices $i$ with $x_i \not = x_1$. Then $a_{ij} = 0$ for $i \in I$
and $j \not \in I$. Thus $I = \{1, \ldots, n\}$ by condition (2) and
$x$ is a multiple of $m = (1, \ldots, 1)$.
\end{proof}

\end{document}
